% © 2026 Ing. David Jaroš — CC BY-NC-ND 4.0
%
% This work is licensed under a Creative Commons Attribution-NonCommercial-NoDerivatives
% 4.0 International License (CC BY-NC-ND 4.0).
%
% License History: Earlier drafts (up to v0.3) were released under CC BY 4.0.
% From v0.4 onward, all material is released under CC BY-NC-ND 4.0 to protect
% the integrity of the theoretical work during ongoing academic development.
%
% See LICENSE.md for full license text.

\ifdefined\INCLUDEMODE
\else
  \documentclass[12pt,a4paper]{article}
  \usepackage[a4paper, margin=2.5cm]{geometry}
  \usepackage{amsmath,amssymb,amsthm}
  \usepackage{booktabs}
  \usepackage{tcolorbox}
  \newtheorem{conjecture}{Conjecture}[section]
  \newtcolorbox{statusbox}[1][]{colback=gray!8!white,colframe=black!60,
    arc=3pt,boxrule=1pt,fonttitle=\bfseries,title=Status Summary,#1}
  \title{Self-Consistency Residuum vs One-Loop QED Correction}
  \author{Ing.\ David Jaro\v{s}}
  \date{2026-05-10}
  \begin{document}
  \maketitle
\fi

\section{Numerical inputs and outputs}

Using
\begin{equation}
f(p)=\frac{5p-1}{p+1}-\ln p,
\end{equation}
we get
\begin{equation}
f(137)=0.036541.
\end{equation}

Observed offset (used only as post-check, not as derivation input):
\begin{equation}
\Delta\alpha^{-1}_{\mathrm{exp}}=137.035999-137=0.035999,
\end{equation}
with ratio
\begin{equation}
\frac{f(137)}{\Delta\alpha^{-1}_{\mathrm{exp}}}=1.0151
\quad(\text{difference }1.505\%).
\end{equation}

\section{One-loop leptonic QED estimate at $\Lambda=M_{\mathrm{Pl}}$}

Using
\begin{equation}
\Delta_{\mathrm{QED},\ell}=\frac{\alpha}{3\pi}\ln\!\left(\frac{\Lambda}{m_\ell}\right),
\end{equation}
with $\alpha=1/137.036$ and $\Lambda=M_{\mathrm{Pl}}=1.22\times10^{19}\,$GeV,
\begin{align}
\Delta_e &= 0.045245,\\
\Delta_\mu &= 0.035768,\\
\Delta_\tau &= 0.033583,\\
\sum\Delta_{\mathrm{QED}} &= 0.114595.
\end{align}
Hence
\begin{equation}
\frac{f(137)}{\sum\Delta_{\mathrm{QED}}}=0.3189,
\end{equation}
so mismatch is $68.11\%$.

Therefore the identification
$\varepsilon_{\mathrm{UBT}}\equiv f(137)=\varepsilon_{\mathrm{QED}}$
in this one-loop leptonic implementation is \textbf{not supported}; classification [NC].

\section{UV scale solving}

Solving $\Delta_{\mathrm{QED}}(\Lambda)=f(137)$:
\begin{itemize}
  \item electron-only formula gives
  $\Lambda\approx1.58\times 10^{14}\,$GeV ($\Lambda/M_{\mathrm{Pl}}\approx1.30\times10^{-5}$),
  \item three-lepton summed formula gives
  $\Lambda\approx3.10\times10^{4}\,$GeV ($\Lambda/M_{\mathrm{Pl}}\approx2.54\times10^{-15}$).
\end{itemize}
Neither equals the Planck scale in this setup.

\begin{conjecture}[Residual link remains open]
The close numerical proximity $f(137)\approx\Delta\alpha^{-1}_{\mathrm{exp}}$
may indicate a structural relation, but the present one-loop leptonic
$\Delta_{\mathrm{QED}}$ model with $\Lambda=M_{\mathrm{Pl}}$ does not reproduce
$f(137)$ and cannot be used as proof.
\end{conjecture}

\begin{statusbox}
$f(137)=0.036541$\\
$\Delta\alpha^{-1}_{\mathrm{exp}}=0.035999$\\
Shoda (vs observed offset): 98.495\% (ratio 1.0151)\\
$\sum\Delta_{\mathrm{QED}}$ (lepton loops, $\Lambda=M_{\mathrm{Pl}}$): 0.114595\\
Interpretace: NC (mismatch 68.11\% against one-loop leptonic $\sum\Delta_{\mathrm{QED}}$)\\
UV škála $\Lambda$ pro přesnou shodu: $1.58\times10^{14}$ GeV (electron-only) / $3.10\times10^4$ GeV (three-lepton sum), not Planck\\
Klasifikace: NC / OPEN for deeper mechanism
\end{statusbox}

\ifdefined\INCLUDEMODE
\else
  \end{document}
\fi
